\section{Method}
\label{sec:method}

\begin{figure}[!t]
\centering
\includegraphics[width=\columnwidth]{figures/fig_task_scenario.png}
\caption{Task scenario: both orbits observe the slope before the event; rapid response may have only the target post-event orbit. Paired-orbit training supports target-only deployment.}
\label{fig:scenario}
\end{figure}

\begin{figure*}[!t]
\centering
\includegraphics[width=\textwidth]{figures/fig_architecture_current.pdf}
\caption{Complete Teacher--Student paths. Counter context guides policy; values remain target-side. TGD is off in CG0. Both distillation retrievals use Student target memory. The Teacher encoder drawings share weights.}
\label{fig:architecture}
\end{figure*}

\subsection{Problem Formulation}
For each paired training sample, $X_t$ and $X_c$ denote target- and counter-orbit pre/post-event SAR stacks and $Y$ is the landslide mask. The Counter-Guided Search Teacher receives $(X_t,X_c)$ during training; the Target-Only Search Student receives $X_t$ during training and inference. The role of $X_c$ is to teach an evidence-search operation executable from $X_t$. Geometry codes are used to define diagnostic evaluation regions. The primary Student does not receive geometry as an input or use it to weight its loss.

\subsection{Counter-Guided Target Evidence Search Teacher}
An encoder produces multi-scale target and counter features,
\begin{equation}
T=E(X_t), \qquad C=E(X_c).
\label{eq:teacher_features}
\end{equation}
At each reference location $p$, the Teacher predicts a joint policy over $L=3$ FPN levels (P2--P4), with $K=4$ candidates per level: one fixed center and three learned points. Weights are normalized across all twelve candidates. Offsets in feature-pixel units are bounded by $2\tanh(\cdot)$:
\begin{equation}
\begin{aligned}
\pi_T(p)&=\{\Delta p^T_{p,l,k},a^T_{p,l,k}\}_{l,k}=\mathcal{P}(T,C)_p,\\
&\hspace{20mm}\sum_{l,k} a^T_{p,l,k}=1.
\end{aligned}
\label{eq:teacher_policy}
\end{equation}
The search operator reads target features only:
\begin{equation}
\begin{aligned}
R_T(p)&=\mathcal{S}(T;\pi_T)_p\\
&=\sum_{l=1}^{L}\sum_{k=1}^{K}a^T_{p,l,k}W_v^{l}T_l(p_l+\Delta p^T_{p,l,k}),
\end{aligned}
\label{eq:teacher_search}
\end{equation}
where $p_l$ maps the query to level $l$, fractional locations use bilinear interpolation, and $W_v^l$ is a learned value projection. A decoder maps $R_T$ to $P_T$. The counter orbit guides policy generation, but no counter feature enters the searched value path or decoder.

\subsection{Target-Only Student and Search Distillation}
The Student forms target features, policy, and searched representation as
\begin{equation}
S=E_S(X_t), \qquad \pi_S=\mathcal{P}_S(S), \qquad R_S=\mathcal{S}(S;\pi_S).
\label{eq:student_search}
\end{equation}
It receives two transfer signals. Policy distillation teaches \emph{where} to search by aligning teacher and student offsets and weights, $\widetilde{\mathcal{L}}_{\mathrm{policy}}(\pi_S,\pi_T)$. Evidence distillation teaches \emph{what} to retrieve by applying both policies to the \emph{same student target memory}:
\begin{equation}
R_{T\rightarrow S}=\mathcal{S}(S;\pi_T), \qquad
\widetilde{\mathcal{L}}_{\mathrm{evidence}}=\ell(R_S,R_{T\rightarrow S}).
\label{eq:shared_memory}
\end{equation}
Thus, the Student neither reconstructs counter-orbit features nor receives a hidden fusion path; it learns how the privileged Teacher policy should be realized on target-side information available at deployment.

\subsection{Learning Objectives}
The primary Teacher uses pixelwise binary cross-entropy (BCE) segmentation loss without geometry-specific weighting. The selected Student first learns from target-only segmentation and subsequently adds policy and evidence alignment. Its objective is
\begin{equation}
\mathcal{L}_S=\mathcal{L}_{\mathrm{BCE}}+
\lambda_p\mathcal{L}_{\mathrm{policy}}+
\lambda_e\mathcal{L}_{\mathrm{evidence}},
\label{eq:student_loss}
\end{equation}
where policy alignment uses the Jensen--Shannon distance between sampling distributions and evidence alignment uses cosine distance between representations retrieved from the same Student target memory. Five initial epochs optimize segmentation alone. The fifth epoch also calibrates the distillation magnitudes without applying their gradients; the calibrated components are added from epoch six with a cosine decay. Teacher outputs, Student target memory, and backbone features used to recompute the distilled Student policy are detached. Consequently, the distillation terms update the context/policy branch, while segmentation trains the complete Student. This stop-gradient construction does not directly distill the backbone or value projection.
